\documentclass[11pt]{article}
\usepackage{../shared/luxcover}
\usepackage{amsmath,amssymb}
\usepackage{hyperref}
\usepackage{booktabs}
\usepackage{enumitem}
\usepackage{xcolor}

\title{Lux Security Legal Framework}
\author{Zach Kelling\\\textit{Lux Industries, Inc.}\\\texttt{research@lux.network}}
\date{May 2023}

\begin{document}

\luxcoverpage

\begin{abstract}
[Security Tokenization Framework

A unique vehicle to tokenize any asset or risk associated with receivables
\end{abstract}

\tableofcontents
\newpage


A unique vehicle to tokenize any asset or risk associated with receivables


Lux Partners Limited offers security tokenization services under a robust regulatory framework both domestically and internationally. These services are compliant with the applicable legal provisions, drawing upon Luxembourg's securitisation act of 22 March 2004 and the United States' securities regulations.


In the United States, Lux follows the S-1 filing process with the Securities and Exchange Commission (SEC). This registration enables Lux to market security tokens to both accredited and non-accredited investors, broadening the investor base and allowing a wider array of individuals and institutions to participate in these tokenization opportunities.


The regulatory framework allows the tokenization of various types of assets, risks, revenues, and activities, making security tokenization accessible to all types of investors, whether institutional or individual. Issuers will use security tokens as an alternative to traditional bank funding.


A vast array of assets can be tokenized, including securities (stocks, loans, subordinated or non-subordinated bonds), risks associated with debt (commercial and others), tangible and intangible property, and any activity with a certain value or future income.


Security Tokenization Process


Security tokenization undertakings established by Lux Partners Limited are incorporated to tokenize any type of assets or risks associated with receivables or any activities conducted by third parties. The tokenization process involves acquiring risks from an originator by issuing a security token, the value of which, along with associated yields, are linked to the underlying asset.


Security tokenization is a process where investors buy security tokens, i.e., transfer cash to a security tokenization vehicle (STV), with the aim of obtaining the returns from the investments made by that vehicle. STVs are assets that produce a predictable cash flow or grant the right to a future cash flow, transforming these assets into security tokens. Investors, however, also bear the risk associated with any uncertainty in these cash flows.


These security tokens can be qualified as Asset-Backed Tokens (ABT), because the underlying assets serve as collateral for the investment. Therefore, the investor generally carries two risks: the uncertainty of a future cash flow and the risk of valuation of the underlying asset.


For instance, a bank may decide to tokenize the risks associated with their real-estate loan portfolio, effectively removing this risk from their balance sheet. The purchaser of these tokens is entitled to the cash flow related to the interest paid by property owners and to the underlying property in case of default by the owners.


Leading the Way in Security Tokenization Services


Lux Partners Limited is committed to safeguarding and advancing the interests of clients by providing services that help them transform their assets and financial futures. Our services are characterized by the following:


In-depth understanding of all the key regulatory and developmental issues


Commitment to finding the most suitable solutions for establishing and managing security tokenization vehicles


Strong and open relationships with clients, based on clear communication and trust


Independence


Lux Partners Limited offers a variety of services for clients, including administration, tax advice, and security token issuance. For further details about our security tokenization services, please visit our dedicated website: lux.fund.]


\end{document}
